\section{Introduction}\label{sec:intro}

Individual heterogeneity is ubiquitous in economic data. Consumers differ in their sensitivity to prices, workers respond differently to job training, and patients vary in their treatment effects. Empirical researchers have long recognized this heterogeneity and built it into their models, from discrete choice models \citep{mcfadden1974conditional, train2009discrete} to personalized pricing \citep{dube2023personalized} and heterogeneous treatment effects more broadly. But estimating heterogeneous parameters at the individual level, and doing valid statistical inference on the resulting objects, has remained a difficult task. Traditional approaches typically assume a parametric distribution for heterogeneity (e.g., normal random coefficients) or restrict attention to a small number of discrete types. \citet{farrell2021deep, farrell2025deep} develop a framework in which deep neural networks directly output individual-level structural parameters $\theta(x)$, and influence function corrections deliver valid confidence intervals despite the regularization bias inherent in neural network estimation. The discrete choice model in our application is one example; the framework applies to any structural model with a differentiable loss.

In the spirit of \citet{nevo2000practitioner}, this paper is a practitioner's guide to the FLM framework. We walk through the entire pipeline --- from data preparation to model specification to inference to counterfactuals --- using a running example that illustrates every step. We provide complete code, a companion Python package (\texttt{deep-inference}), and a simulation study validating the method. Related work on deep learning for causal inference includes DragonNet \citep{shi2019dragonnet}, which adapts network architecture for treatment effect estimation; GDR-Learners \citep{melnychuk2026gdr}, which extend doubly robust methods to heterogeneous effects; and \citet{hetzenecker2024deep}, who apply similar ideas to discrete choice models.

\subsection{The Idea in Brief}

Consider a structural model with individual-level parameters. A researcher observes outcomes $Y_i$, treatments or product characteristics $T_i$, and individual covariates $X_i$ for $i = 1, \ldots, n$ observations. The structural model specifies a loss function $\ell(Y_i, T_i, \theta)$, and the individual's true parameters $\theta^*(X_i)$ minimize the expected loss conditional on their type $X_i$.

The key insight of FLM is: let a neural network directly output the structural parameters. Train a network $\hat{\theta}(X_i)$ by minimizing the structural loss across observations. Then use influence functions to correct for the regularization bias that neural networks inevitably introduce.

The framework applies broadly. For example, in a \emph{logit} demand model, $\ell$ is the binary cross-entropy and $\theta(X_i) = (\alpha_i, \beta_i)$ captures individual-level intercept and price sensitivity. In a \emph{Poisson} count model for store visits, $\ell$ is the Poisson negative log-likelihood and $\beta_i$ is the individual's response to a promotional campaign. In a \emph{multinomial choice} model, $\ell$ is the softmax cross-entropy and $\theta(X_i)$ is a vector of taste parameters across product attributes.
The target $\mu^* = E[H(X, \theta^*(X), \tilde{t})]$ can be an average parameter ($H = \theta_k$, e.g., mean price sensitivity), a marginal effect ($H = g'(\theta'\tilde{t}) \cdot \beta$), an economic quantity like expected revenue ($H = \tilde{t} \cdot g(\theta'\tilde{t})$), or a treatment effect contrast ($H = g(\theta'\tilde{t}_1) - g(\theta'\tilde{t}_0)$, i.e., a difference-in-differences-style estimand).

Why is correction necessary? Neural networks use regularization (weight decay, dropout, early stopping) to prevent overfitting. This regularization shrinks estimates toward zero, making naive standard errors, computed from the Hessian of the structural loss, systematically too small. The resulting confidence intervals severely undercover. In our simulations, naive coverage can fall below 20\% when the nominal level is 95\%.

The influence function (IF) correction accounts for the estimation noise in $\hat{\theta}(X)$. The corrected estimator for a target functional $\mu^* = E[H(X, \theta^*(X), \tilde{t})]$ is:
\begin{equation}\label{eq:psi_intro}
    \hat{\psi}_i = H_i - H_{\theta,i} \cdot \hat{\Lambda}_i^{-1} \cdot \ell_{\theta,i}
\end{equation}
where $H_i$ is the target evaluated at the estimated parameters, $H_{\theta,i}$ is the Jacobian of the target with respect to $\theta$, $\hat{\Lambda}_i$ is the expected Hessian of the loss, and $\ell_{\theta,i}$ is the score (gradient of the loss). The average $\bar{\psi} = n^{-1}\sum_i \hat{\psi}_i$ is the debiased estimator, and its standard error provides valid inference.

The influence function has a long history in robust statistics \citep{hampel1974influence} and semiparametric efficiency theory \citep{newey1994large}. In this setting, it serves a specific purpose: it corrects for the first-order bias introduced by estimating $\theta$ with a regularized neural network. The key property is \emph{Neyman orthogonality}: the expectation of $\psi$ is insensitive to first-order perturbations in $\hat{\theta}$ around $\theta^*$. Concretely, $\partial_\theta E[\psi(Y, T, \theta, \Lambda)]|_{\theta^*} = 0$. This means that even though $\hat{\theta}$ converges to $\theta^*$ at a slow rate (typical for neural networks), the debiased estimator $\bar{\psi}$ achieves the parametric $\sqrt{n}$ rate. The correction term $H_\theta \Lambda^{-1} \ell_\theta$ is constructed precisely so that this orthogonality condition holds: it cancels the leading-order bias from the plug-in estimator $\bar{H}$.

\subsection{Running Example: H\&M Fashion Demand}

Our running application uses the H\&M Personalized Fashion Recommendations dataset \citep{hm2022kaggle}, a publicly available dataset of 1.4 million fashion transactions. We estimate a multinomial choice model for online dress purchases with heterogeneous price sensitivity across consumers.

The application demonstrates the full pipeline. First, \emph{representation learning}: we train two-tower contrastive embeddings \citep{oord2018representation} to learn dense representations of consumers and products from purchase histories and visual features, where consumer embeddings capture latent taste dimensions and item embeddings capture visual and categorical product attributes. Second, \emph{choice model estimation}: consumer embeddings become the covariates $X$ that drive heterogeneity, and the neural network maps each consumer's 64-dimensional embedding to their structural parameters $\theta_{\text{price}}(X_i)$ and $\theta_{\text{style},k}(X_i)$. Third, \emph{influence function inference}: the \texttt{deep-inference} package handles cross-fitting, Lambda estimation, influence function assembly, and confidence interval computation, with the entire pipeline running automatically from a five-line custom loss function. Finally, \emph{counterfactual analysis}: we compute average price elasticities with valid confidence intervals and simulate the revenue impact of a 10\% price increase, with uncertainty bands that properly account for the neural network estimation step.

\subsection{What This Guide Covers}

Section~\ref{sec:framework} develops the framework: the structural model, target functionals, why naive inference fails, the IF correction, and the three Lambda estimation regimes. Section~\ref{sec:estimation} covers estimation in practice: data requirements, the algorithm step by step, cross-fitting, Lambda estimation, network architecture, diagnostics, and common pitfalls. Section~\ref{sec:examples} provides three self-contained worked examples of increasing complexity (linear, logit, and custom Poisson). Section~\ref{sec:application} presents the full H\&M application with results. Section~\ref{sec:simulations} validates the pipeline through simulations, including semi-synthetic studies with high-dimensional covariates. Section~\ref{sec:conclusion} discusses extensions and other structural models. Appendix~\ref{app:implementation} provides annotated code and an installation guide.

Throughout, we emphasize the \emph{practitioner's} perspective: what choices must be made, what can go wrong, and what to look for in the output.
